本研究は、交通分野における量子技術開発と社会的要請の関係を評価するため、交通・物流分野の社会的要求段階と既存・想定量子最適化技術の技術段階を対応づけ、幅・深さを中間証拠としてそのギャップを評価する枠組みを提示した。

分析では、既存文献の幅、深さ、ゲート数、ショット数、実行時間、実行可能性、資源推定を抽出し、それらをQRL証拠として分類した。
さらに、社会側要求を問題、データ・ワークフロー、制約、資源、技術段階に分け、SRL-QRL対応表とギャップ表を構成した。
表Dでは現状の規模、資源、ハードウェア・ワークフロー、実行可能性、幅・深さのトレードオフギャップを整理し、本文の状態遷移要約表では6つの技術社会改善シナリオに対する評価結果を示した。

結果として、技術側だけの改善や資源証拠だけの改善では、多くの場合、実行可能性、ワークフロー、社会側準備度の判定条件が残ることが示された。
また、社会側要求だけが高まる場合には、量子側証拠が追いつかず、ギャップは拡大する。
実用候補に近づくのは、資源、実行可能性、ワークフロー、社会側準備度が同時に改善される技術社会の均衡改善シナリオに限られる。
したがって、EV・物流システムの設計者・運用者・評価者が量子最適化技術をPoC、資源推定、または実用候補として検討できる条件は、量子回路資源だけでなく、符号化、実行可能性の処理、ハイブリッドワークフロー、社会側運用条件の共同整備として評価する必要がある。

本研究の限界は、準備度の状態遷移評価が性能予測ではなく、文献ベースの順序尺度スコアリングと判定条件に基づく構造評価である点にある。
今後は、TMS・配車ワークフローや古典ベースラインに関する一次資料を追加し、状態遷移評価で示した次に必要な証拠を実験・実証計画へ接続する必要がある。
また、全30行のシナリオ結果をもとに、ダッシュボード化や対象領域を変えた比較へ拡張することも可能である。
